\chapter{Anal Itching (Pruritus Ani)}

\section*{Definition}
Anal itching (pruritus ani) is itching or irritation of the skin around the anus. It may be short-lived or persistent and often has more than one contributing factor. Moisture, stool leakage, irritants, skin disease, infection, and anorectal conditions can all trigger an itch-scratch cycle.

\section*{When to See a Doctor}
\begin{itemize}
\item Itching that does not improve with gentle skin care or keeps returning
\item Itching accompanied by rectal bleeding, discharge, a lump, significant pain, or swelling
\item Visible worms or intense nighttime itching, especially in a child
\item Skin changes such as redness, fissures, or lichenification
\item Nocturnal itching (common with pinworms)
\item A history of diabetes, immunosuppression, inflammatory bowel disease, or a possible sexually transmitted infection
\end{itemize}

\section*{How It Develops}
Pruritus ani involves irritation of the perianal skin, which triggers an itch-scratch cycle. Scratching causes further skin damage and inflammation. Moisture, stool, mucus, and irritants can damage the skin barrier. Scratching can cause fissures and secondary infection, but fungal or bacterial infection should not be assumed without compatible findings.

\section*{Causes}
\begin{itemize}
\item Moisture or irritation from stool leakage, diarrhea, incontinence, sweating, soap residue, vigorous wiping, or excessive cleaning
\item Fecal soiling/incontinence leading to perianal moisture
\item Hemorrhoids or anal fissures
\item Pinworm infection (Enterobius vermicularis), especially in children
\item Fungal, bacterial, parasitic, or viral infection when supported by examination
\item Contact dermatitis from soaps, lotions, toilet paper, or laundry detergents
\item Psoriasis or seborrheic dermatitis involving the perianal area
\item Food or drink triggers in some people; these are individual and should not be assumed
\item Anal fistula or other anorectal conditions
\end{itemize}

\section*{Related Symptoms}
\begin{itemize}
\item Hemorrhoids
\item Anal fissure
\item Rectal discharge
\item Rectal pain
\item Perianal rash
\item Diarrhea or fecal incontinence
\item Itching or rash elsewhere, or vulvar itching when a nearby skin or vaginal condition is involved
\end{itemize}

\section*{Medical Evaluation}
\begin{itemize}
\item Visual inspection of the perianal area
\item A morning adhesive-tape test or other appropriate test for pinworm when the history and examination suggest it; routine stool testing is not the usual test for pinworm
\item Swabs, skin scrapings, or other tests only when infection is suspected
\item Anoscopy or other colorectal evaluation when bleeding, discharge, a mass, bowel-change symptoms, or another anorectal condition is suspected
\item Patch testing or dermatology referral when allergic contact dermatitis is suspected
\end{itemize}

\section*{Treatment Options}
\begin{itemize}
\item Treat the identified cause, such as pinworm infection, dermatitis, hemorrhoids, fissure, or infection; do not use antifungals or antibiotics without evidence of infection
\item A short course of weak topical corticosteroid may reduce inflammation when advised, but prolonged use can thin and irritate perianal skin
\item A barrier cream such as zinc oxide or petroleum jelly can protect irritated skin; topical antihistamines may themselves cause contact dermatitis
\item Behavioral modification to break the itch-scratch cycle
\item Sitz baths with plain warm water after bowel movements
\item Dietary changes only when a clear personal trigger is identified; avoid unnecessarily restrictive diets
\item Referral to a dermatologist or colorectal specialist for refractory cases
\end{itemize}

\section*{Self-Care and Home Management}
\begin{itemize}
\item Clean the perianal area gently with plain warm water after bowel movements
\item Pat dry (do not rub) and avoid harsh soaps or fragranced wipes
\item Apply a thin layer of petroleum jelly or zinc oxide paste as a barrier
\item Wear loose-fitting, breathable cotton underwear
\item Avoid scratching; keep fingernails short
\item Avoid a food or drink only if it clearly worsens symptoms; routine elimination of caffeine or spicy foods is not necessary for everyone
\item Ask a pharmacist or clinician before using hydrocortisone near the anus, and use it sparingly for a short period only; stop if irritation worsens
\end{itemize}

\section*{References}
\begin{thebibliography}{99}
\bibitem{anal_itching:pruritusani1} Chaudhry V, Qazi MK, Ahmed Z. ``Pruritus ani: a review of the literature.'' Am J Gastroenterol. 2020;115(9):1387-1394.
\bibitem{anal_itching:pruritusani2} Stern E, Prout T. ``Pruritus ani.'' In: StatPearls. StatPearls Publishing; 2023.
\bibitem{anal_itching:pruritusani3} Markell KW, Billingham RP. ``Pruritus ani: etiology and management.'' Surg Clin North Am. 2010;90(1):125-135.
\bibitem{anal_itching:pruritusani4} Bowyer A, McColl I. ``The treatment of pruritus ani.'' J R Soc Med. 1996;89(2):95-97.
\bibitem{anal_itching:pruritusani5} Harrington CI, Lewis JW. ``Pruritus ani.'' BMJ. 2016;353:i2562.
\bibitem{anal_itching:pruritusani6} MacLean J, Russell D. ``Pruritus ani.'' Can Fam Physician. 2010;56(6):543-546.
\bibitem{anal_itching:pruritusani7} American Society of Colon and Rectal Surgeons. Pruritus ani: expanded patient information. Updated 2025.
\end{thebibliography}
